\section{Background and Related Work}
\label{sec:related-work}

Most PUBG prediction work uses post-match aggregate statistics to predict final
placement or rank \cite{kagglePubgFinishPlacement,sen2021pubgFinalRank,lee2025pubgRank}.
That framing is useful for final-outcome modeling, but it hides when risk
changes inside a match and how the shrinking safe zone changes the meaning of
movement or combat.

The closest prior work is Sagong et al.'s phase-based PUBG survival framework
\cite{sagong2025hunterCamper}. They use phase-specific aggressiveness and
activeness scores, called FightScore and ActiveScore, to classify players into
strategy types and predict next-phase survival. Their main contribution for our
purpose is not a particular classifier, but the framing: PUBG survival should
be analyzed by phase.

We extend that framing in two ways. First, ecology has long studied organisms
whose habitats shrink or fragment: habitat-fragmentation and movement-ecology
research treats viable space, movement, exposure, and density as inherently
spatial and temporal quantities
\cite{fahrig2003fragmentation,nathan2008movement}. The shrinking safe zone
imposes the same kind of contraction on teams, so we borrow this lens as an
analytic analogy, not a biological claim, and it points to the same
conclusion: survival conditions should be modeled spatio-temporally rather
than as aggregate scores. Second, we represent
team-to-team context explicitly. Each phase becomes a graph whose edges encode
proximity, combat, and shared safe-zone pressure
\cite{sanchezLengeling2021gnnIntro}. Thus the contribution is not merely using
a graph model. It is to reconstruct official telemetry as a spatio-temporal and
relational representation, cast survival as a discrete-time hazard in the
event-history tradition \cite{allison1982discrete,singer1993discrete}, and
test next-phase elimination risk against strong non-graph baselines.
